\documentclass[12pt,a4paper]{article}
\usepackage[utf8]{inputenc}
\usepackage[T1]{fontenc}
\usepackage[brazil]{babel}
\usepackage{amsmath, amssymb}
\usepackage{graphicx}
\usepackage{booktabs}
\usepackage{enumitem}
\usepackage{hyperref}
\usepackage{geometry}
\usepackage{caption}
\usepackage{tikz}
\usetikzlibrary{arrows.meta, shapes.geometric, shapes.misc, positioning, fit}

\geometry{left=3cm,right=2cm,top=3cm,bottom=2cm}

\title{\textbf{Metodologia de Preparação, Síntese e Validação de Dados (QIP)}}
\author{Trabalho de Conclusão de Curso}
\date{\today}

\begin{document}
\maketitle

\section{Visão geral}
A metodologia adotada se baseia em: (i) preparar a base real; (ii) calibrar uma heurística de decisão multi-classe; (iii) sintetizar dados condicionais por classe com um cWGAN\textendash GP; (iv) treinar, com os dados gerados na etapa anterior, um modelo heurístico baseado em \emph{softmax}; e (v) validar o desempenho em um conjunto externo não utilizado em nenhuma etapa anterior. A abordagem foi concebida para cenários de amostras escassas e classes desbalanceadas, garantindo rastreabilidade e separação rigorosa entre \emph{dados para análise} e \emph{dados para validação}.

\section{Preparação dos dados e ``Seleção de Dados''}
\subsection{Base original e limpeza}
A base original, extraída do questionário QIP, contém variáveis numéricas e um alvo categórico. Amostras com rótulo \emph{desconhecido} (e.g., \texttt{não/nao}) são removidas antes de qualquer etapa posterior. Esse procedimento evita viés no condicionamento do gerador e impede contaminação do conjunto de validação.\footnote{Implementação de remoção e normalização do alvo nos scripts de pré-processamento.}

\subsection{Seleção de atributos orientada}
Adota-se o critério de \textbf{não utilizar variáveis demográficas} e/ou outras perguntas cuja influência seja difícil de analisar por conhecimentos heurísticos. Essa filtragem reduz dimensionalidade e preserva apenas atributos com contribuição informativa prévia.\footnote{O script registra as colunas mantidas e excluídas. Opcionalmente, pode-se incorporar regras adicionais (por exemplo, \emph{regra de ouro} com base em aba de pontuações específicas).}

\subsection{Tratamento de classes raras}
Classes com menos de um limiar mínimo de amostras são removidas (p.ex., $\leq 2$ respostas). Essa decisão evita falhas em procedimentos que requerem pelo menos duas amostras por classe (e.g., vizinhanças do SMOTE) e estabiliza a estimação de parâmetros condicionais.\footnote{Contagens por classe, limiar adotado e remoções são registradas em log e planilhas de saída.}

\section{Heurística de decisão (calibração de $W$)}
\subsection{Mudança de faixa e motivação}
Inicialmente, a heurística considerou pesos $W\in[0,1]$. Após testes, adotou-se $W\in[-1,1]$, o que promoveu melhor discriminação: ao reforçar a classe correta, penaliza-se as demais. Observou-se incremento de desempenho em validações internas (e.g., $\approx36\%\rightarrow45\%$).

\subsection{Formulação}
Seja $X\in\mathbb{R}^{n\times d}$ o conjunto de atributos e $W\in\mathbb{R}^{(d+1)\times K}$ os pesos (incluindo viés). As probabilidades são $P=\mathrm{softmax}(X_b W)$, com $X_b$ denotando $X$ acrescido de 1s (viés). A calibração maximiza uma métrica alvo (macro top-$k$) sob restrições: máscara de atributos permitidos (consistente com a seleção anterior), projeção em $[-1,1]$ nos pesos dos atributos e \emph{early stopping} por platô.

\section{Síntese condicional com cWGAN-GP}
\subsection{Separação analítica vs. validação externa}
A base é dividida em duas metades aproximadamente estratificadas: \textbf{(A)} metade para \emph{análise} (treino/ajuste) e \textbf{(B)} metade para \emph{validação externa}, que não participa de nenhuma etapa de ajuste (nem cWGAN, nem heurística).

\subsection{Balanceamento e condicionamento}
Na metade A, a distribuição por classe é balanceada com \emph{oversampling} (SMOTE com \emph{fallback} para amostragem aleatória com reposição, quando necessário). Em seguida, treina-se um cWGAN-GP condicional por classe, operando em espaço escalado $[0,1]$ e revertendo ao espaço original para análise. Variáveis binárias e inteiras (0..10) são arredondadas na exportação dos dados sintéticos.

\subsection{Geração de dados sintéticos}
O cWGAN é aplicado exclusivamente na metade A para gerar amostras virtuais, ampliando a cobertura de padrões sob as distribuições condicionais de classe. A metade B permanece preservada para validação externa.

\section{Treinamento final (softmax/heurística) com dados gerados}
O modelo heurístico final é treinado sobre o conjunto ampliado da metade A (reais + sintéticos do cWGAN), mantendo as mesmas restrições de projeto (máscara de atributos, projeção $[-1,1]$, \emph{early stopping}). Essa etapa consolida um classificador \emph{softmax} com melhor generalização em presença de classes desbalanceadas.

\section{Validação externa}
A heurística final é aplicada \textbf{exclusivamente} na metade B (jamais vista), reportando-se: (i) macro top-$k$ (para $k=1,2,3$); (ii) métricas por classe (acerto/suporte); (iii) gráficos/relatórios de diagnóstico do processo de calibração.

\section{Métricas auxiliares (qualidade do sintético)}
Durante o treino do cWGAN e na amostragem final, computam-se indicadores auxiliares:
\begin{itemize}[leftmargin=1.5em]
    \item \textbf{TSTR} (Train on Synthetic, Test on Real): macro-F1 de um classificador simples treinado em sintético e testado no real (metade A);
    \item \textbf{KS mediano} por atributo (real vs. sintético) no espaço original;
    \item \textbf{Gap de correlação} (norma de Frobenius) entre matrizes de correlação real e sintética;
    \item \textbf{Diferença média/máxima} de proporções por classe (real vs. sintético).
\end{itemize}
Esses sinais complementam, mas não substituem, a validação externa na metade B.

\section{Boas práticas e salvamento de resultados}
Todos os artefatos são salvos em planilhas com abas para: dados reais de treino (\texttt{DATA\_TRAIN}), sintéticos (\texttt{DATA\_SYNTH}), validação externa (\texttt{DATA\_HOLDOUT}), bem como relatórios de seleção de atributos, parâmetros de normalização e histórico de avaliação.

\section{Fluxograma do processo}
\begin{figure}[h!]
\centering
\begin{tikzpicture}[
    node distance=7mm and 14mm,
    every node/.style={font=\small},
    >={Latex},
    box/.style={draw, rounded corners=2pt, align=center, minimum width=40mm, minimum height=8mm, fill=gray!5},
    io/.style={draw, trapezium, trapezium left angle=70, trapezium right angle=110, align=center, minimum width=40mm, minimum height=8mm, fill=blue!5},
    decision/.style={draw, diamond, aspect=2, align=center, inner sep=1pt, fill=orange!10},
    dashedbox/.style={draw, dashed, rounded corners=2pt, align=center, minimum width=40mm, minimum height=8mm, fill=gray!3}
]

% Entrada
\node[io] (start) {TDados (real)};
\node[box, below=of start] (clean) {Remoção de rótulos desconhecidos\\ (\texttt{não/nao})};
\node[box, below=of clean] (sel) {Seleção de atributos\\ (excluir demográficos e itens difíceis)};
\node[box, below=of sel] (rare) {Remoção de classes raras\\ (suporte $\leq$ limiar)};

% Split
\node[decision, below=of rare] (split) {Divisão estratificada 50/50};

% Metade A (Análise)
\node[box, below left=14mm and 18mm of split] (A) {Metade A (análise)};
\node[box, below=of A] (bal) {Balanceamento (SMOTE/ROS)};
\node[box, below=of bal] (gan) {Treino cWGAN-GP\\ (condicionado por classe)};
\node[box, below=of gan] (aug) {Geração de sintéticos\\ (apenas A)};
\node[box, below=of aug] (heur) {Treinamento final\\ (softmax/heurística $W\in[-1,1]$)};

% Metade B (Validação)
\node[box, below right=14mm and 18mm of split] (B) {Metade B (validação externa)};
\node[box, below=of B] (hold) {Mantida intocada};
\node[box, below=of hold] (test) {Aplicar heurística final\\ e medir desempenho};

% Ligações principais
\draw[->] (start) -- (clean);
\draw[->] (clean) -- (sel);
\draw[->] (sel) -- (rare);
\draw[->] (rare) -- (split);
\draw[->] (split) -- node[above, sloped]{A} (A);
\draw[->] (split) -- node[above, sloped]{B} (B);
\draw[->] (A) -- (bal);
\draw[->] (bal) -- (gan);
\draw[->] (gan) -- (aug);
\draw[->] (aug) -- (heur);
\draw[->] (B) -- (hold);
\draw[->] (hold) -- (test);

% Caixa de métricas/diagnóstico
\node[box, below=24mm of heur] (metrics) {Métricas auxiliares (TSTR, KS, correlação, dif. de classe)};

% Setas para métricas
\draw[->] (heur) -- (metrics);
\draw[->] (test) -- ++(0,-6mm) -| (metrics);

% Grupo visual (fit) para as metades
\node[draw, rounded corners=3pt, fit=(A)(bal)(gan)(aug)(heur), inner sep=4pt, label=above:{\footnotesize Pipeline da Metade A}] {};
\node[draw, rounded corners=3pt, fit=(B)(hold)(test), inner sep=4pt, label=above:{\footnotesize Pipeline da Metade B}] {};

\end{tikzpicture}
\caption{Fluxograma consolidado: preparação, síntese, treinamento final e validação externa.}
\end{figure}

\section{Parâmetros e escolhas de projeto}
\begin{itemize}[leftmargin=1.5em]
    \item \textbf{Seleção de dados}: exclusão explícita de variáveis demográficas/itens de difícil interpretação; limiar mínimo de suporte por classe;
    \item \textbf{Heurística}: pesos de atributos projetados em $[-1,1]$; viés livre; \emph{early stopping};
    \item \textbf{cWGAN-GP}: dimensão de ruído $z$, largura/profundidade das camadas, $n_{\text{critic}}$, penalidade de gradiente, ruído de instância opcional;
    \item \textbf{Balanceamento}: alvo total de amostras e estratégia (SMOTE com \emph{fallback} ROS/duas passadas);
    \item \textbf{Avaliação}: frequência de avaliação, janelas de média móvel do \emph{gap}, critérios de parada por estabilidade/platô, baterias TSTR/KS/correlação.
\end{itemize}

\section{Considerações éticas e de privacidade}
Para reduzir riscos de reidentificação, avaliam-se distâncias de vizinhos reais-sintéticos e aplicam-se limiares de descarte quando necessário, mantendo a divisão estrita entre \emph{análise} e \emph{validação externa}.

\end{document}